\documentclass[10pt,letterpaper]{article}
\usepackage{cornell-notes}

\title{Clause 05.13.01: Kinds Of Literals}
\author{}
\date{}
\setDocTitle{Clause 05.13.01: Kinds Of Literals}
\setDocSubtitle{C++ 2024}
\setDocOwner{C++ 2024 Cornell Notes}
\setDocDate{\today}
\setCornellCollection{C++ 2024 Cornell Notes}
\setCornellUnitType{Clause}
\setCornellUnitNumber{05.13.01}
\setCornellUnitTitle{Kinds Of Literals}
\hypersetup{pdftitle={Clause 05.13.01: Kinds Of Literals},pdfsubject={Cornell notes},pdfauthor={}}

\begin{document}
\maketitle
\begin{CornellOverviewBox}{Overview}
\textbf{Classification:} Core language - Normative\\[2pt]
\textbf{Learning objective:} Understand the rules, terminology, and semantic consequences associated with Kinds of literals.
\end{CornellOverviewBox}\begin{CornellSummaryBox}{Summary}
{\sffamily\bfseries\color{Accent} Summary}\par\vspace{3pt}Section 5.13.1 focuses on Kinds of literals. Use the listed grammar productions together with the semantic constraints and disambiguation rules referenced by the section. When applying it, preserve the distinction between mandatory requirements, recommendations, permissions, and behavior deliberately left undefined, unspecified, or implementation-defined.
\end{CornellSummaryBox}

\end{document}
